\chapter{Implementation, Reproducibility, and Final Reporting Pipeline}

\section{Software and Repository Stack}

The thesis stack is implemented inside a single repository that contains the CYCLES simulator interface, CyclesGym environments, experiment runners, reporting utilities, documentation, and the LaTeX thesis workspace. This co-location keeps code, frozen artifacts, and generated writing assets under one versioned project. The active training entrypoints remain \texttt{experiments/fertilization/train.py} and \texttt{experiments/crop\_planning/train.py}, while \texttt{run\_experiments\_7\_3\_2026.py} preserves the design of the full 113-run matrix that became the final campaign.

What changed for the release-ready version of the thesis is the reporting layer. Instead of treating ad hoc runner CSVs as the primary evidence surface, the repository now freezes the completed campaign in \texttt{artifacts/final\_successful\_runs/final\_113/} and rebuilds all thesis assets from the canonical reporting outputs generated under \texttt{final\_113/reporting/}. This design sharply reduces the risk of result drift between code, artifacts, and prose.

\section{Data Inputs and Provenance}

The thesis depends on three categories of inputs. The first is environmental input: the Pakistan weather file and Pakistan soil file wired into the current training and evaluation defaults. The second is economic input: yearly crop and nutrient price series generated from official-source-derived pipelines. The third is operational input: the frozen experiment bundles, recovered rerun exports, and the canonical reporting summaries that standardize the final evidence set.

The localized price series is especially important because it converts cost awareness from a narrative claim into executable reward logic. Reward terms that depend on nutrient cost are no longer interpreted through a generic legacy profile alone; they are tied to yearly Pakistan-oriented series assembled from FAOSTAT and NFDC-linked sources \citep{faostatprices2026,nfdcprices2026}.

\input{tables/generated/artifact_inventory.tex}

\begin{figure}[H]
    \centering
    \includegraphics[width=0.92\textwidth]{data_provenance_diagram.pdf}
    \caption{Data provenance flow from primary sources to final thesis evidence. The figure clarifies how external source documents become executable weather, calendar, and pricing inputs and then enter the canonical reporting layer.}
    \label{fig:data-provenance}
\end{figure}

\input{tables/generated/data_source_provenance.tex}

\section{Training and Evaluation Pipeline}

The training pipeline used by the thesis stack is summarized in Figure~\ref{fig:training-pipeline}. Command-line configuration determines which environment is built, how vectorized wrappers are applied, which algorithm is trained, and where evaluation and summary artifacts are written.

\begin{figure}[H]
    \centering
    \includegraphics[width=0.90\textwidth]{training_pipeline_diagram.pdf}
    \caption{Training and evaluation pipeline used by the thesis experiments. The same pipeline is responsible for the summary JSON outputs later canonicalized into the final reporting dataset.}
    \label{fig:training-pipeline}
\end{figure}

Standardized per-run summary JSON files are central to this design. They preserve the quantities that the thesis actually claims over, including deterministic return, stochastic return, Pakistan holdout return for fertilization, runtime metadata, and provenance fields that indicate whether a row came from the original curated bundle set or a recovered rerun.

\input{tables/generated/summary_output_schema.tex}

\section{Canonical Reporting Pipeline}

The reporting pipeline is equally important because the thesis does not rely only on aggregate returns. Hierarchical environments can emit weekly nutrient logs, yearly crop decisions, and compliance summaries where these directories are present. More importantly, the final release now contains a repo-level canonical report builder, \texttt{scripts/build\_final\_reports.py}, which resolves the frozen bundle set into one authoritative reporting directory containing run-level metrics, grouped metrics, statistical tests, an artifact audit, and a summary JSON for top-line counts.

\begin{figure}[H]
    \centering
    \includegraphics[width=0.90\textwidth]{reporting_pipeline_diagram.pdf}
    \caption{Canonical reporting workflow from frozen artifacts to thesis assets. The important change is that the LaTeX tables and figures now consume \texttt{final\_113/reporting/} instead of provisional runner summaries.}
    \label{fig:reporting-pipeline}
\end{figure}

This pipeline also makes the March 14, 2026 guarded hierarchical reruns explicit. Those reruns are treated as authoritative for the hierarchical branch, but they are preserved with provenance as corrected guarded reruns rather than being silently merged into the original March 7--11 raw hierarchical narrative.

\section{Observed Runtime and Final Matrix Completion}

The completed campaign now has an observed runtime profile rather than only a planned one. Figure~\ref{fig:final-runtime} and Table~\ref{tab:runtime-summary} summarize the actual runtime footprint of the frozen evidence set, including the longer guarded hierarchical reruns.

\begin{figure}[H]
    \centering
    \includegraphics[width=0.84\textwidth]{final_runtime_profile.pdf}
    \caption{Observed runtime profile of the frozen final evidence set. This is based on completed runs rather than on a historical planning estimate.}
    \label{fig:final-runtime}
\end{figure}

\input{tables/generated/runtime_summary.tex}

Completion status is now equally concrete. The final matrix contains 113 frozen rows, including 16 recovered replacements, 12 guarded hierarchical reruns, and 4 recovered DQN reruns. These counts are summarized in Figure~\ref{fig:final-matrix-status} and Table~\ref{tab:current-results-status}.

\begin{figure}[H]
    \centering
    \includegraphics[width=0.84\textwidth]{final_matrix_status.pdf}
    \caption{Composition of the frozen final evidence set after canonicalization.}
    \label{fig:final-matrix-status}
\end{figure}

\input{tables/generated/current_results_status.tex}

\section{Reproducibility Controls}

The repository now includes several reproducibility controls that are directly relevant to thesis quality.
\begin{itemize}[leftmargin=1.2cm]
    \item The final evidence base is frozen in a single artifact tree, \texttt{final\_113}, rather than being reconstructed from multiple competing summary folders.
    \item Canonical reporting is generated by one builder that standardizes run-level, grouped, statistical, and audit outputs.
    \item The thesis workspace rebuilds figures and tables from canonical reporting outputs rather than from manual transcription.
    \item Recovery provenance is explicit: replacement rows, corrected guarded hierarchical reruns, and descriptive-only DQN rows are all carried forward with labels instead of being hidden.
    \item Administrative thesis metadata remains isolated in a dedicated file so that document formatting changes do not disturb the technical chapters.
\end{itemize}

\input{tables/generated/verification_checks.tex}

Table~\ref{tab:verification-checks} is intentionally concrete. It records the release verification steps that were executed for the final cleaned repository. The goal is not to claim formal software verification. The goal is to show that the most thesis-critical pathways, namely canonical report generation, thesis asset regeneration, full PDF rebuild, count cross-checks, and stale-text removal, were all checked directly.

\section{Why This Chapter Matters for the Final Thesis}

A thesis with strong motivation and attractive figures can still struggle under committee scrutiny if the software pipeline is not reproducible or if different folders tell different stories. This chapter exists to prevent that risk. It explains how localized inputs, frozen experiment bundles, canonical reporting outputs, and rebuildable LaTeX assets fit together into one release-ready pipeline that supports the final results chapter without manual patchwork.
